% Generated from data/qjapanimation.yaml via scripts/generate_qjapanimation_tex.py.
% Edit the YAML data or this template, not qjapanimation.tex directly.
\documentclass[10pt,a4paper,landscape]{ltjsarticle}

\usepackage[
  top=10mm,
  bottom=12mm,
  left=10mm,
  right=10mm
]{geometry}
\usepackage{array}
\usepackage{longtable}
\usepackage[table]{xcolor}
\usepackage{hyperref}
\usepackage{enumitem}

\hypersetup{
  colorlinks=true,
  urlcolor=blue!55!black,
  linkcolor=black,
  pdftitle={量子情報に関連する話題が登場するアニメ・漫画（ときどきその他）一覧},
  pdfauthor={生田力三}
}

\setlength{\parindent}{0pt}
\setlength{\LTpre}{0pt}
\setlength{\LTpost}{4mm}
\setlength{\tabcolsep}{1.5mm}
\renewcommand{\arraystretch}{1.18}
\arrayrulecolor{black!25}

\newcolumntype{Y}[1]{>{\raggedright\arraybackslash}p{#1}}
\newcolumntype{C}[1]{>{\centering\arraybackslash}p{#1}}

\newcommand{\QJATableHeader}{%
  \hline
  \rowcolor{black!6}
  \textbf{年} & \textbf{タイトル} & \textbf{内容} & \textbf{備考} \\
  \hline
}

\begin{document}

\begin{center}
  {\Large\bfseries 量子情報に関連する話題が登場するアニメ・漫画（ときどきその他）一覧\par}
  \vspace{1.5mm}
  生田力三\par
  \vspace{0.7mm}
  {\small\color{black!55}最終更新: 2026/09/02\par}
\end{center}
\vspace{1mm}

このページでは、量子情報に関連する話題が登場するアニメ・漫画などをまとめています。量子力学そのものに加えて、量子通信、量子コンピュータ、量子もつれ、量子テレポーテーション、量子暗号、量子アルゴリズム、多世界解釈、シュレーディンガーの猫など、量子情報の周辺に出てくるキーワードや設定を見つけるたび、気まぐれに記録しています。

\vspace{1mm}
\begin{itemize}[leftmargin=5mm,itemsep=1mm,topsep=1mm]
  \item 2015年頃から少しずつ更新している、発展途上の一覧です。
  \item 漫画・アニメなど複数のメディアに作品がある場合、特に断りがない限りアニメの年を記載しています。
  \item 関連情報をご存じでしたら、お知らせいただけると幸いです。
  \item シュレーディンガーの猫が登場する作品は非常に多いため、特別なものを除き後半の表にまとめています。
  \item \href{https://doi.org/10.18910/92424}{「量子の未来」をめぐる23の話題}（2023年8月17日）の「話題19」で本ページが紹介されました。
\end{itemize}

\section*{量子情報・量子力学に関する話題}
\small
\begin{longtable}{C{13mm}|Y{46mm}|Y{106mm}|Y{99.5mm}}
\QJATableHeader
\endfirsthead
\QJATableHeader
\endhead
  1996 & 機動戦艦ナデシコ & ボソンジャンプ &  \\
  \hline
  2001 & 千と千尋の神隠し & 千尋の服が北野先生の「量子力学の基礎」(2010)の表紙と同じ(？)デザイン & 意図的？ \\
  \hline
  2003 & 狂四郎2030 & 光明が量子コンピュータを開発 & 15巻付近 \\
  \hline
  2005 & ノエイン & 量子力学、多世界論 & 作品の感想はとてもこの場では言えない \\
  \hline
  2006 & 涼宮ハルヒの憂鬱 & 高い指向性を持つ不可視帯域のコヒーレント光 & 22話 溜息III 長門 \\
  \hline
  2006 & ゼーガペイン & エンタングルメント、量子サーバ、量子エラー訂正、など & 基本的に量子情報の話ばかり \\
  \hline
  2006 & マブラヴオルタネイティヴ & 因果律的量子論 & 物理の話が何時間か延々と続く \\
  \hline
  2007 & 電脳コイル & 「量子回路のある特殊な基盤パターンが過去例を見ないほど高性能なアンテナになる事に…(中略)…回路が電磁波以外に何かを受信していたのを発見したのよ。人間の意識よ。」 & 23話 原理は分からなかったが回路のコピーは簡単だったとか \\
  \hline
  2007 & 機動戦士ガンダム00 & 量子コンピュータ「ヴェーダ」 &  \\
  \hline
  2007 & 涼宮ハルヒの分裂 & 「粒子のような振る舞いをしながら波動としての仕事もする、まるで光のように」 & 古泉 \\
  \hline
  2007 & えれくとりっく・えんじぇぅ & 「まるで量子の風みたいに」 & ミクだしとりあえず入れとくかみたいな、ヤスオP、唄：初音ミク \\
  \hline
  2009 & サマーウォーズ & Shor のアルゴリズム & 本編の感想はとてもここには書けない \\
  \hline
  2009 & クォンタム・ファミリーズ & 量子回路、量子コンピュータ、量子脳計算機科学、ほか、香ばしい言葉がたくさん & 東浩紀著。第23回三島由紀夫賞。アニメになってもそんなにおかしくはない。ネットワーク感が前面に出ていて良い。 \\
  \hline
  2010 & ムダヅモ無き改革 & 「そうさ どちらも正しくてどちらも同時に存在する!! シュレーディンガーの猫のケツもクソまみれだぜ」 & 4巻の湯川ヒデキ \\
  \hline
  2011 & 問題児たちが異世界から来るそうですよ？ & 立体交差並行宇宙 & この手のを入れるときりがないが一応 \\
  \hline
  2011 & 涼宮ハルヒの直観 & 「仮に並行世界があったとして、それは無限に存在するものじゃないのかい。こう、エヴェレット的に」 & 佐々木 \\
  \hline
  2011 & 量子世界のあなたへ -- Quantum Lovers -- & シュレーディンガーの猫、並行世界、ヒルベルト空間、波動関数のΦ、経路積分、不確定性原理、など & キーワードてんこ盛り、あいんしゅたいんP、唄：初音ミク \\
  \hline
  2011 & STEINS;GATE & 世界線、タイムマシン、並行世界 & 一応 \\
  \hline
  2012 & ヨルムンガンド & 量子コンピュータ &  \\
  \hline
  2012 & DCIII & 多世界解釈 & 一応 \\
  \hline
  2012 & アクセルワールド & 量子回路、脳内の光量子回路 &  \\
  \hline
  2012 & 僕 & 観測するまで事象は決まらない & きづきあきら+サトウナンキ \\
  \hline
  2012 & Quantum Entanglement & 量子のレベルで、干渉し合う波動、もつれて絡まっていく、など & 作詞作曲：かめりあ、唄：初音ミク \\
  \hline
  2013 & 蒼き鋼のアルペジオ & 「量子通信、良好」 & 良好とか気安く言っててつらい \\
  \hline
  2014 & キャプテンアース & エンタングルリンク、エンタングルゲート &  \\
  \hline
  2014 & 俺、ツインテールになります。 & 「私の脳内の無限長(中略)にエンタングルメント」 & 10話母のセリフ \\
  \hline
  2014 & クロスアンジュ & 並行世界、ふたつの地球、統一理論を旋律に、量子フィールド &  \\
  \hline
  2014 & 失われた未来を求めて & シュレーディンガーの猫、多世界解釈、余剰次元、量子チューリングマシン、チューリングテスト、量子サイバネティクス、量子エンタングルメント開始 & 5話、9話 \\
  \hline
  2015 & 放課後のプレアデス & どちらの存在でもない重ね合わせの状態 & 10話 \\
  \hline
  2015 & 電波教師 & 質量をもつ物体の瞬間移動についての実証実験、最終目的は、通称どこでもドアの製造 & 22話 素粒子物理学研究所CERM所長の会見 \\
  \hline
  2015 & 乱歩奇譚 & 量子論とカオスの融合がどうたらこうたら。ナミコシの計算ノートにケット記法を用いた謎の計算あり。 &  \\
  \hline
  2015 & アルドノアゼロ & 「量子テレポートを利用して自分のコピーを作っているのかもしれないその時コピーの元原稿を壊さず残していることで分身しているんだ」 & 意味不明 \\
  \hline
  2015 & 美男高校地球防衛部LOVE! & 「この個体は死んでもあるが生きてもあるよ」「なんだそのシュレーディンガー的な言い訳は」 & 第2話 \\
  \hline
  2015 & iショウジョ+ & P≠NP問題 & 漫画1巻v13 \\
  \hline
  2016 & クロムクロ & 「彼の脳にはナノマシンとEPR相関関係にある回路があるのでしょうが(後略)」 & 26話 \\
  \hline
  2016 & ビッグオーダー & パラレルワールド、量子的にもつれた心や魂 & ９話 \\
  \hline
  2016 & アンジュヴィエルジュ & センサーが未来から敵のデータを量子通信で過去に送る & 2話 \\
  \hline
  2017 & 亜人ちゃんは語りたい & 「人間が観測することによって状態が変化するものがあるのだ、電子一粒がどうのという世界における量子と呼ばれるものだ」 & 10話 \\
  \hline
  2017 & Chaos;Child & 「ディラックの海」「た、たしか量子力学の中の概念ですよね」 & 6話 \\
  \hline
  2017 & ひなろじ & 「ベルのパラドックスレベルを検出して行方をおっていますが(後略)」 & 3話 \\
  \hline
  2017 & ID-0 & IDO「砂どころか素粒子レベルまで…トンネル効果ですり抜けて…」 & 5話 \\
  \hline
  2017 & 十二大戦 & 未「儂が戦士として持つ器量は量子力学トンネル効果を自在に操り物質を透過する技能ですじゃ」 & 5話 \\
   &  & ハンドレッドクリック、シュレーディンガーの猫 & 11話 \\
  \hline
  2017 & 景の海のアペイリア & 二重スリット実験 &  \\
  \hline
  2018 & BEATLESS & 「レイシア級--ヒューマノイドインターフェースエレメンツ--(中略)量子コンピュータを搭載したデバイスを装備し、ネットワーク支援なしで高度な判断を行う能力をもつ」 & 1話 \\
  \hline
  2018 & 青春ブタ野郎はバニーガール先輩の夢を見ない & シュレーディンガーの猫、ラプラスの悪魔、量子もつれ、量子テレポーテーション、観測云々 & 胸踊るキーワードがたくさん \\
  \hline
  2018 & されど罪人は竜と踊る & 「プランク定数hを操作し、森羅万象を生み出す力、咒力」 & TVアニメ \\
  \hline
  2019 & ルパン三世 グッバイパートナー & (スーパー)量子コンピューター & 他にもタイムクリスタル、対称性の破れ、AIやらシンギュラリティやら雑ワード多め \\
  \hline
  2019 & HELLO WORLD & 量子記憶装置アルタラ & それはさておき聖地巡礼したい \\
  \hline
  2019 & ぬるぺた & ショアのアルゴリズムによる素因数分解に関する計算式 & 3話、黒板 \\
  \hline
  2019 & 青春ブタ野郎はゆめみる少女の夢を見ない & 量子もつれとか観測理論とか & 難解すぎる \\
  \hline
  2020 & エチカの時間 & 量子コンピュータのランダムサーキットサンプリングに挑戦 & 13話、世界中の暗号通貨も解読してました \\
  \hline
  2020 & 神様になった日 & チップ型量子コンピュータ & アツい、常温だけに() \\
  \hline
  2021 & シン・エヴァンゲリオン劇場版 & 量子テレポート & すごくよかったです \\
  \hline
  2021 & EX-ARM & 量子コンピュータ & 漫画第1巻は2015年 \\
  \hline
  2022 & アカデミックSOS & 「スパコンも量子コンピュータもあって使いたい放題らしいです」 & 1話、闇すぎて読んでるとなぜか動悸が…… \\
  \hline
  2022 & code：ノストラ & 未来を予測できる量子コンピュータ「ノストラ」「ファティマ」 & 出だし読むとキョドッてしまうので注意が必要 \\
  \hline
  2022 & 異世界薬局 & 「任意の次元の量子の情報は〜」 & TVアニメ \\
  \hline
  2023 & Q.E.D. iff ―証明終了― & 量子コンピュータ、量子ゲート & 第25巻「第四のゲート」 \\
  \hline
  2024 & Electricity & 量子もつれがテーマだとかなんとか & 宇多田ヒカルの曲、アニメ映画の主題歌とかになっててもおかしくない、知らんけど \\
  \hline
  2024 & アンダーニンジャ & 「量子コンピュータでも解読困難な最新のポスト量子暗号(PQC)」 & 第116話 \\
  \hline
  2024 & 涼宮の劇場 & 量子コンピュータ、量子情報、量子力学的重ね合わせ、エンタングルメント、光子のみで構成された準天体規模の量子演算処理システム、ホログラフィック原理、など楽しげなワードがたくさん & 量子論に対する人類の理解がまだ足りてないことを長門に言わせた上で、長門本人にごちゃつかせないで古泉に語らせることで不正確なことでも自然に受け入れられる説明になってて神(早口) \\
  \hline
  2024 & 地球防衛隊X & 量子AI「Akashic down」 & 31話、外生体との共同開発とのこと \\
  \hline
  2024 & 科学×冒険サバイバル！ & 量子力学の研究者「クオ博士」 & ふつうに身近な存在である \\
  \hline
  2025 & 青春ブタ野郎はサンタクロースの夢を見ない & 量子もつれ、ヒルベルト空間 & ヒルベルト空間はタイトルのみ \\
  \hline
\end{longtable}
\normalsize

\clearpage
\section*{シュレーディンガーの猫}
\small
\begin{longtable}{C{13mm}|Y{46mm}|Y{106mm}|Y{99.5mm}}
\QJATableHeader
\endfirsthead
\QJATableHeader
\endhead
  -- & 絶望先生 & シュレーディンガーの猫 &  \\
  \hline
  -- & 咲-Saki- & シュレーディンガーの猫 &  \\
  \hline
  -- & とある科学の超電磁砲 & シュレーディンガーの猫 &  \\
  \hline
  -- & うみねこのなく頃に & シュレーディンガーの猫 &  \\
  \hline
  -- & 神様のメモ帳 & シュレーディンガーの猫 &  \\
  \hline
  -- & デュラララ!! & シュレーディンガーの猫 &  \\
  \hline
  2009 & 夏のあらし！ & シュレーディンガーの猫、観測によって値が確定 &  \\
  \hline
  2014 & 中二病でも恋がしたい! & シュレーディンガーの猫 & 2期1話 \\
  \hline
  2014 & 異能バトルは日常系のなかで & シュレーディンガーの猫 &  \\
  \hline
  2015 & プリズマイリヤツヴァイヘルツ！ & シュレーディンガーの猫 & 4話アバンタイトル \\
  \hline
  2016 & Dimension W & シュレーディンガーの猫 & 11話 \\
  \hline
  2016 & ディバインゲート & シュレーディンガーの猫 & 4話 \\
  \hline
  2021 & 出会って5秒でバトル & シュレーディンガーの猫 & 152② \\
  \hline
\end{longtable}
\normalsize

\vfill
{\footnotesize\color{black!55}年は作品公開年を基本とし、一部に年不明の項目を含みます。}

\end{document}
